% =====================================================================
\chapter{Cuestiones}
\label{cap:test}
% =====================================================================

Veinte cuestiones de opción múltiple, de 2 puntos cada una (40 puntos en total). Hay al menos dos cuestiones de cada módulo de la fase. Hay una sola respuesta correcta en cada cuestión y las respuestas erróneas no restan. La clave y su justificación están en el \cref{ap:soluciones}.

\begin{cuestiones}
% ----------------------------------------------------------- Módulo 0.1
\item \pt{2} La negación del enunciado «$\forall\varepsilon>0\ \exists N\in\N\ \forall n\ge N:\ \abs{a_n-L}<\varepsilon$» es:
\begin{opciones}
  \item $\exists\varepsilon>0\ \forall N\in\N\ \exists n\ge N:\ \abs{a_n-L}\ge\varepsilon$;
  \item $\forall\varepsilon>0\ \exists N\in\N\ \forall n\ge N:\ \abs{a_n-L}\ge\varepsilon$;
  \item $\exists\varepsilon>0\ \forall N\in\N\ \forall n\ge N:\ \abs{a_n-L}\ge\varepsilon$;
  \item $\exists\varepsilon\le0\ \exists N\in\N\ \forall n\ge N:\ \abs{a_n-L}\ge\varepsilon$.
\end{opciones}

\item \pt{2} Sea $n\in\Z$. El contrarrecíproco de «si $n^{2}$ es par, entonces $n$ es par» es:
\begin{opciones}
  \item si $n$ es par, entonces $n^{2}$ es par;
  \item si $n^{2}$ es impar, entonces $n$ es impar;
  \item si $n$ es impar, entonces $n^{2}$ es impar;
  \item $n^{2}$ es par y $n$ es impar.
\end{opciones}

% ----------------------------------------------------------- Módulo 0.2
\item \pt{2} La función $f\colon\Z\to\Z$, $f(n)=2n+1$, es:
\begin{opciones}
  \item biyectiva;
  \item sobreyectiva, pero no inyectiva;
  \item ni inyectiva ni sobreyectiva;
  \item inyectiva, pero no sobreyectiva.
\end{opciones}

\item \pt{2} Sean $f\colon A\to B$ y $g\colon B\to C$. Si $g\circ f$ es inyectiva, entonces necesariamente:
\begin{opciones}
  \item $f$ es inyectiva;
  \item $g$ es inyectiva;
  \item $f$ es sobreyectiva;
  \item $f$ y $g$ son inyectivas.
\end{opciones}

% ----------------------------------------------------------- Módulo 0.3
\item \pt{2} El valor de $\displaystyle\sum_{k=1}^{20}(2k-1)$ es:
\begin{opciones}
  \item 380
  \item 400
  \item 420
  \item 441
\end{opciones}

\item \pt{2} Si $x,y>0$, $\log_2 x=5$ y $\log_2 y=3$, entonces $\log_2\bigl(x^{2}/y\bigr)$ vale:
\begin{opciones}
  \item $10/3$
  \item $2$
  \item $7$
  \item $13$
\end{opciones}

\item \pt{2} La representación de $-5$ en complemento a dos con 8 bits es:
\begin{opciones}
  \item \texttt{10000101}
  \item \texttt{11111010}
  \item \texttt{00000101}
  \item \texttt{11111011}
\end{opciones}

% ----------------------------------------------------------- Módulo 0.4
\item \pt{2} El valor exacto de $\cos 75^{\circ}$ es:
\begin{opciones}
  \item $\dfrac{\sqrt6+\sqrt2}{4}$
  \item $\dfrac{\sqrt6-\sqrt2}{4}$
  \item $\dfrac{\sqrt3-1}{2}$
  \item $\dfrac{\sqrt2-\sqrt6}{4}$
\end{opciones}

\item \pt{2} Para $n\ge2$, la suma de las $n$ raíces $n$-ésimas de la unidad es:
\begin{opciones}
  \item $1$
  \item $n$
  \item $0$
  \item $-1$
\end{opciones}

% ----------------------------------------------------------- Módulo 0.5
\item \pt{2} El límite $\displaystyle\lim_{x\to0}\frac{1-\cos x}{x^{2}}$ vale:
\begin{opciones}
  \item $1/2$
  \item $0$
  \item $1$
  \item no existe.
\end{opciones}

\item \pt{2} La derivada de la sigmoide $\sigma(x)=1/(1+e^{-x})$ es:
\begin{opciones}
  \item $\sigma(x)^{2}$
  \item $1-\sigma(x)^{2}$
  \item $e^{-x}\sigma(x)$
  \item $\sigma(x)\bigl(1-\sigma(x)\bigr)$
\end{opciones}

\item \pt{2} El polinomio de Maclaurin de grado~3 de $f(x)=\sin x$ es:
\begin{opciones}
  \item $x-\dfrac{x^{3}}{3}$
  \item $x-\dfrac{x^{3}}{6}$
  \item $x+\dfrac{x^{3}}{6}$
  \item $1-\dfrac{x^{2}}{2}$
\end{opciones}

% ----------------------------------------------------------- Módulo 0.6
\item \pt{2} El valor de $\displaystyle\int_{0}^{1}x\,e^{x^{2}}\,\dd x$ es:
\begin{opciones}
  \item $e-1$
  \item $e/2$
  \item $(e-1)/2$
  \item $(e+1)/2$
\end{opciones}

\item \pt{2} ¿Cuál de las siguientes series es convergente?
\begin{opciones}
  \item $\displaystyle\sum_{n\ge1}\frac{1}{\sqrt n}$
  \item $\displaystyle\sum_{n\ge1}\frac{1}{n^{1{,}01}}$
  \item $\displaystyle\sum_{n\ge1}\frac{1}{n}$
  \item $\displaystyle\sum_{n\ge1}\frac{n}{n+1}$
\end{opciones}

% ----------------------------------------------------------- Módulo 0.7
\item \pt{2} El gradiente de $f(x,y)=x^{2}y+e^{y}$ en el punto $(1,0)$ es:
\begin{opciones}
  \item $(0,2)$
  \item $(2,1)$
  \item $(0,1)$
  \item $(1,2)$
\end{opciones}

\item \pt{2} Una función de clase $C^{2}$ tiene un punto crítico en el que su matriz hessiana es $\begin{pmatrix}2&3\\3&2\end{pmatrix}$. Ese punto es:
\begin{opciones}
  \item un mínimo local;
  \item un máximo local;
  \item un punto de silla;
  \item un punto en el que el criterio del hessiano no decide.
\end{opciones}

% ----------------------------------------------------------- Módulo 0.8
\item \pt{2} El número de palabras distintas (con o sin sentido) que se obtienen reordenando las letras de BANANA es:
\begin{opciones}
  \item 720
  \item 120
  \item 36
  \item 60
\end{opciones}

\item \pt{2} Una enfermedad afecta al 1\,\% de la población. Una prueba da positivo en el 99\,\% de los enfermos y en el 5\,\% de los sanos. La probabilidad de estar enfermo si la prueba da positivo es, aproximadamente:
\begin{opciones}
  \item $0{,}99$
  \item $0{,}95$
  \item $0{,}17$
  \item $0{,}01$
\end{opciones}

% ----------------------------------------------------------- Módulo 0.9
\item \pt{2} Se toma una muestra aleatoria simple de tamaño $n=5$ de una población con varianza $\sigma^{2}=10$. La esperanza del estimador $\widehat\sigma^{2}=\frac1n\sum_{i=1}^{n}(X_i-\Xbar)^{2}$ es:
\begin{opciones}
  \item $10$
  \item $12{,}5$
  \item $2$
  \item $8$
\end{opciones}

\item \pt{2} Con una muestra de $n=100$ datos de una población normal con $\sigma=10$ conocida se obtiene $\xbar=50$. El intervalo de confianza del 95\,\% para la media es:
\begin{opciones}
  \item $[48{,}36;\ 51{,}65]$
  \item $[48{,}04;\ 51{,}96]$
  \item $[30{,}40;\ 69{,}60]$
  \item $[49{,}80;\ 50{,}20]$
\end{opciones}
\end{cuestiones}
